\documentclass[aspectratio=169]{beamer}

\usetheme{Madrid}
\usecolortheme{default}

\usepackage[T1]{fontenc}
\usepackage[utf8]{inputenc}
\usepackage[french]{babel}
\usepackage{lmodern}
\usepackage{graphicx}
\usepackage{booktabs}

\title[Titre court]{Titre de la presentation}
\subtitle{Sous-titre optionnel}
\author{Prenom Nom}
\institute{Institution}
\date{\today}

\begin{document}

\begin{frame}
  \titlepage
\end{frame}

\begin{frame}{Plan}
  \tableofcontents
\end{frame}

\section{Contexte}

\begin{frame}{Motivation}
  \begin{itemize}
    \item Presenter le probleme.
    \item Expliquer pourquoi il est important.
    \item Annoncer l'approche proposee.
  \end{itemize}
\end{frame}

\section{Methode}

\begin{frame}{Approche}
  \begin{block}{Idee principale}
    Decrire la methode en une ou deux phrases.
  \end{block}

  \begin{enumerate}
    \item Etape 1.
    \item Etape 2.
    \item Etape 3.
  \end{enumerate}
\end{frame}

\section{Resultats}

\begin{frame}{Resultats principaux}
  \begin{table}
    \centering
    \begin{tabular}{lcc}
      \toprule
      Methode & Score & Temps \\
      \midrule
      Reference & 0.80 & 12 s \\
      Proposee & 0.91 & 9 s \\
      \bottomrule
    \end{tabular}
  \end{table}
\end{frame}

\section{Conclusion}

\begin{frame}{Conclusion}
  \begin{itemize}
    \item Resume du message principal.
    \item Limites de l'etude.
    \item Perspectives.
  \end{itemize}
\end{frame}

\begin{frame}
  \centering
  \Large Merci pour votre attention.
\end{frame}

\end{document}

